%% !TeX root=GUATThesis-main.tex
%% !TEX TS-program = XeLaTeX
%بسته های برای تایپ متون ریاضی و محیط های شماره دار
\usepackage{amsthm,amssymb,amsmath,amsfonts}			% Writing math
\usepackage{mathtools}
%% \usepackage{upgreek}  %%برای ایستاده نمایش دادن شکل ایستاده نمادهای یونانی
\usepackage[shortlabels]{enumitem}
%%%%%%%%%%%%%%%%%%%%%%%%%%%%%%%%%%%%%%%%%%%%%%%%%%%
% بسته‌ای برای تنطیم حاشیه‌های بالا، پایین، چپ و راست صفحه
\usepackage[top=40mm, bottom=20mm, left=20mm, right=20mm]{geometry}
%%%%%%%%%%%%%%%%%%%%%%%%%%%%%%%%%%%%%%%%%%%%%%%%%%%
% بسته ای برای رسم انواع مختلف نمودارها
\usepackage{tikz}
%%%%%%%%%%%%%%%%%%%%%%%%%%%%%%%%%%%%%%%%%%%%%%%%%%%
% بسته‌‌ای برای ظاهر شدن شکل‌ها و تصاویر متن
\usepackage{epsf,graphicx}								% Including graphics
\usepackage{caption, subcaption}
\usepackage{xcolor,colortbl}
\definecolor{green}{rgb}{0.5,1,0.5}
\usepackage{makecell}
%%%%%%%%%%%%%%%%%%%%%%%%%%%%%%%%%%%%%%%%%%%%%%%%%%
%دستوری برای وارد کردن واژه‌نامه انگلیسی به فارسی
\newcommand\persiangloss[2]{#1\dotfill\lr{#2}\\}
%\newcommand{\done}{\cellcolor{teal}5- اطلاعات رساله} %{0.9}
%\newcommand{\hcyan}[1]{{\color{teal} #1}}
%%%%%%%%%%%%%%%%%%%%%%%%%%%%%%%%%%%%%%%%%%%%%%%%%%%
\usepackage[stable,bottom]{footmisc}					% Move footnotes to the bottom of page
\usepackage{notoccite}% PREVENTS CITES IN CAPTIONS FROM MISNUMBERING YOUR REFERENCES 
\usepackage{cite}
\usepackage{url}
%%%%%%%%%%%%%%%%%%%%%%%%%%%%%%%%%%%%%%%%%%%%%%%%%%%
% برای مدیریت بر ساختار عناوین، هدر و ... مربوط به فصل، بخش و ...
\usepackage{emptypage}
\usepackage{titlesec}
\usepackage{setspace}									% Change line spacing
%%%%%%%%%%%%%%%%%%%%%%%%%%%%%%%%%%%%%%%%%%%%%%%%%%%
% مورد استفاده در تعاریف با استفاده از newcommand
%\usepackage{xspace}
%%%%%%%%%%%%%%%%%%%%%%%%%%%%%%%%%%%%%%%%%%%%%%%%%%%
%%%% بسته‌ و دستوراتی برای ایجاد لینک‌های رنگی با امکان جهش
\usepackage[colorlinks=true,linkcolor=Blue,urlcolor=Blue,citecolor=Blue]{hyperref}
% چنانچه قصد پرینت گرفتن نوشته خود را دارید، خط بالا را غیرفعال و  از دستور زیر استفاده کنید چون در صورت استفاده از دستور زیر‌‌، 
% لینک‌ها به رنگ سیاه ظاهر خواهند شد و برای پرینت گرفتن، مناسب‌تر خواهد بود.
%\usepackage[pagebackref=false]{hyperref}

%%%%%%%%%%%%%%%%%%%%%%%%%%%%%%%%%%%%%%%%%%%%%%%%%%%
%بسته های مربوط به ایجاد الگوریتم
\usepackage{algorithm,algcompatible}
%%%%%%%%%%%%%%%%%%%%%%%%%%%%%%%%%%%%%%%%%%%%%%%%%%%
%بسته هایی برای مدیریت کپشن در محیط های شناور مانند شکل و جدول و ...
\usepackage[textfont={small},labelfont={rm,small},format=hang,labelsep=quad,justification={centering},aboveskip=2pt,belowskip=4pt]{caption}
\usepackage{subcaption}
%%%%%%%%%%%%%%%%%%%%%%%%%%%%%%%%%%%%%%%%%%%%%%%%%%%

%%%%%%%%%%%%%%%%%%%%%%%%%%%%%%%%%%%%%%%%%%%%%%%%%

%%%%%%%%%%%%%%%%%%%%%%%%%%%%%%%%%%%%%%%%%%%%%%%%%%%
% بسته مورد نیاز برای نوشتن کدهای برنامه نویسی در نوشتار و رنگ
\usepackage{listings}
\usepackage{color}
\definecolor{Blue}{rgb}{0,0,.55}

%%%%%%%%%%%%%%%%%%%%%%%%%%%%%%%%%%%%%%%%%%%%%%%%%%%%%%%
\usepackage{multirow}
\usepackage{array}
\usepackage{pbox}
%تغییر اندازه سطرهای جداول
\renewcommand{\arraystretch}{1.2}
%%%%%%%%%%%%%%%%%%%%%%%%%%%%%%%%%%%%%%%%%%%%%%%%%%%%%%%%%%%%%%%%%%%%%%%%%%%%%%%%%%%%
\usepackage{roundbox,framed}
%%%%%%%%%%%%%%%%%%%%%%%%%%%%%%%%%%%%%%%%%%%%%%%%%%%%%%%%%%%%%%%%%%
%بسته‌ای برای ظاهر شدن «مراجع»  در فهرست مطالب
\usepackage{tocbibind}
\usepackage{zref-perpage}
\zmakeperpage[1]{footnote}
%%%%%%%%%%%%%%%%%%%%%%%%%%%%%%%%
\makeatletter
\renewcommand\tableofcontents{%
	\if@twocolumn
	\@restonecoltrue\onecolumn
	\else
	\@restonecolfalse
	\fi
	\centerline{\huge\bfseries\contentsname
		\@mkboth{%
			\MakeUppercase\contentsname}{\MakeUppercase\contentsname}}%
	\vskip 20\p@
	\@starttoc{toc}%
	\if@restonecol\twocolumn\fi
}
\makeatother
%%%%%%%%%%%%%%%%%%%%%%%%%%%%%%%%
%\usepackage{indentfirst}
% فراخوانی بسته زی‌پرشین و تعریف قلم فارسی و انگلیسی
%\usepackage[extrafootnotefeatures]{xepersian}
\usepackage{xepersian}
%دستور \usepackage{xepersian} باید بعد از همه بسته‌ها (پکیج‌ها)گذاشته شود یعنی زیپرشین آخرین بسته‌ای باید باشد که فراخوانی می‌کنید
%\settextfont[Scale=0.93]{HM XZar}
\settextfont{XB Zar}										% Persian font
\setlatintextfont{Times New Roman}
% تعریف قلم‌های فارسی و انگلیسی برای استفاده در بعضی از قسمت‌های متن
\defpersianfont\titr[Scale=1]{B Titr}
\defpersianfont\Naz[Scale=1]{B Nazanin}
\defpersianfont\traffic[Scale=1]{B Traffic}
% چنانچه می‌خواهید که اعداد داخل فرمول‌ها، فارسی باشد، دستور زیر را فعال کنید
%\setdigitfont{Yas}
%%%%%%%%%%%%%%%%%%%%%%%%%%%%%%%%%%%%%%%%%%

%%%%%%%%%%%%%%%%%%%%%%%%%%%%%%%%%%%%%%%%%

%%%%%%%%%%%%%%%%%%%%%%%%%%%%%%%%%%%%%%%
% Use English digits in equations
\usepackage{amsmath}

%\DefaultMathsDigits

% Default footnotes from left to right
\setfootnoteLR

% Use dash instead of dot in section numbers
%\SepMark{.}										


%% Change fonts and margins of section and subsection titles
%% For chapters please see firstpages.tex
%\titlespacing*{\section}{0pt}{1cm}{0.2cm}
%\titleformat{\section}
%{\fontsize{12}{6}\scshape\bfseries}{\thesection}{1em}{}
%
%\titlespacing*{\subsection}{0pt}{.8cm}{0cm}
%\titleformat{\subsection}
%{\fontsize{11}{6}\scshape\bfseries}{\thesubsection}{1em}{}